\begin{definition}[One fusion layer of the topological projector]%
    \label{def:mpdo_topological_projector_one_fusion}
    \lean{MPOTensor.BNTFusionCoisometryFamily.projectorQBlock}
    \lean{MPOTensor.BNTFusionCoisometryFamily.%
        fusionCoisometry_mul_physTraceTransfer_mul_conjTranspose}
    \lean{MPOTensor.BNTFusionCoisometryFamily.conjugatedProjectorQBlock}
    \leanok
    \uses{def:mpdo_bnt_fusion_coisometry_family}
    Let $P_\gamma$ be operators on the bond space of the vertically read
    tensor $M_\gamma$. Define one fusion layer of the recursive operator $Q$
    by
    \begin{align}
        Q_{\alpha,\beta}
         & =\bigoplus_\gamma
        \chi_{\alpha,\beta,\gamma}\otimes P_\gamma.
        \label{eq:rfp_one_fusion}
    \end{align}
    In the source construction, the terminal matrices are
    $P_\gamma=\tr(M_\gamma)$, with spectral refinement used to obtain
    projections. Here the trace closes the horizontal operator leg and
    leaves the bond indices of the vertically read tensor open.
    Before spectral refinement, the fusion identity implies
    \begin{align}
        U_{\alpha,\beta}\tr(M_\alpha M_\beta)
        U_{\alpha,\beta}^{\dagger}
         & =\bigoplus_\gamma
        \chi_{\alpha,\beta,\gamma}\otimes
        \tr(M_\gamma).
        \label{eq:rfp_terminal_trace}
    \end{align}
    Indeed, summing the forward fusion identity over equal horizontal
    operator indices and distributing matrix multiplication over this finite
    sum gives~(\ref{eq:rfp_terminal_trace}).
    This is one pairwise step in the recursion of
    \cite[lines~999--1010]{Cirac2016MPDO_arXiv}.

    This definition describes one fusion layer. The arbitrary-chain
    recursion, the density decomposition, and the physical commuting Gibbs
    decomposition for $N\geq2$ are given in
    Theorems~\ref{thm:mpdo_topological_projector_recursion},
    \ref{thm:mpdo_topological_density_decomposition}, and
    \ref{thm:mpdo_physical_coordinate_gibbs_at_least_two}, respectively.
    The source boundary at $N=1$ is recorded in
    Remark~\ref{rem:mpdo_topological_gibbs_length_one_boundary}.
\end{definition}

\begin{theorem}[Projection property of one fusion layer]%
    \label{thm:mpdo_topological_projector_one_fusion}
    \lean{MPOTensor.BNTFusionCoisometryFamily.projectorQBlock_eq_unweighted}
    \lean{MPOTensor.BNTFusionCoisometryFamily.projectorQBlock_isStarProjection}
    \lean{MPOTensor.BNTFusionCoisometryFamily.%
        conjugatedProjectorQBlock_isOrthogonalProjection}
    \leanok
    \uses{def:mpdo_topological_projector_one_fusion, thm:mpdo_bnt_length_independent_zipper}
    Let $P_\gamma$ be self-adjoint idempotents on the terminal bond space of
    $M_\gamma$. Suppose that the matrices $\chi_{\alpha,\beta,\gamma}$ are
    diagonal with strictly positive entries, and that the structure
    coefficients
    $c_{\alpha,\beta,\gamma}^{(L)}
        =\tr(\chi_{\alpha,\beta,\gamma}^{L})$ are independent of $L$ for every
    positive length $L$. Then
    \begin{align}
        Q_{\alpha,\beta}
         & =\bigoplus_\gamma
        \Id_{r_{\alpha,\beta}^{\gamma}}\otimes P_\gamma
        \notag
    \end{align}
    is self-adjoint and idempotent, where $r_{\alpha,\beta}^{\gamma}$ is
    the dimension of $\chi_{\alpha,\beta,\gamma}$.
    This is the single recursive fusion step in
    \cite[lines~999--1010]{Cirac2016MPDO_arXiv}.
    Since the fusion map onto the nonzero sectors satisfies
    $U_{\alpha,\beta}U_{\alpha,\beta}^{\dagger}=\Id$, the operator
    $U_{\alpha,\beta}^{\dagger}Q_{\alpha,\beta}U_{\alpha,\beta}$ is also an
    orthogonal projection on the product bond space.
\end{theorem}
\begin{proof}\leanok
    The trace-power hypothesis and length independence force every diagonal
    entry of $\chi_{\alpha,\beta,\gamma}$ to equal one. Thus
    $\chi_{\alpha,\beta,\gamma}=\Id$, while
    $c_{\alpha,\beta,\gamma}^{(L)}=r_{\alpha,\beta}^{\gamma}$ need not equal
    one. Hence
    \begin{align}
        Q_{\alpha,\beta}^{\dagger}
         & =\bigoplus_\gamma(\Id\otimes P_\gamma^{\dagger})
        =Q_{\alpha,\beta},
        \notag                                              \\
        Q_{\alpha,\beta}^{2}
         & =\bigoplus_\gamma(\Id\otimes P_\gamma^{2})
        =Q_{\alpha,\beta}.
        \notag
    \end{align}
    Finally, the coisometry identity
    $U_{\alpha,\beta}U_{\alpha,\beta}^{\dagger}=\Id$ and
    Theorem~\cite[``Coisometric transport of a star projection'']{QICLean2026Blueprint} show that
    $U_{\alpha,\beta}^{\dagger}Q_{\alpha,\beta}U_{\alpha,\beta}$ is an
    orthogonal projection.
\end{proof}

\begin{definition}[Recursive fusion histories and their sequential circuit]%
    \label{def:mpdo_topological_projector_recursion}
    \lean{MPOTensor.BNTFusionCoisometryFamily.FusionHistoryData}
    \lean{MPOTensor.BNTFusionCoisometryFamily.fusionChainTensor}
    \lean{MPOTensor.BNTFusionCoisometryFamily.fusionHistoryWeight}
    \lean{MPOTensor.BNTFusionCoisometryFamily.recursiveProjectorQ}
    \lean{MPOTensor.BNTFusionCoisometryFamily.sequentialFusionCoisometry}
    \leanok
    \uses{def:mpdo_bnt_fusion_coisometry_family}
    Fix an initial label $\alpha$ and append labels successively along a
    finite chain. For total chain length $N\geq1$, the appended-label list
    has length $N-1$; the symbols $W_N$ and $Q_N$ below refer to this total
    chain length. A fusion history records, at every step, the preceding
    final label, the new final label, and an index in the corresponding
    multiplicity space. If $h$ ends at $\gamma(h)$, let
    \begin{align}
        w(h)
         & =\prod_r
        (\chi_{\gamma_{r-1},\beta_r,\gamma_r})_{k_r,k_r}.
        \label{eq:rfp_history_weight}
    \end{align}
    For terminal bond operators $P_\gamma$, define
    \begin{align}
        Q_N(P)
         & =\bigoplus_h w(h)P_{\gamma(h)}.
        \label{eq:rfp_recursive_q}
    \end{align}
    The sequential map $W_N$ is obtained by tensoring the preceding
    map with the identity on the newly appended bond and then applying the
    pairwise fusion coisometry onto the nonzero sectors in each
    preceding-history block.
    This gives the recursive operator and sequential circuit in
    \cite[lines~999--1010]{Cirac2016MPDO_arXiv}.
\end{definition}

\begin{theorem}[Fixed-label recursive fusion identity]%
    \label{thm:mpdo_topological_projector_recursion}
    \lean{MPOTensor.BNTFusionCoisometryFamily.%
        sequentialFusionCoisometry_mul_conjTranspose}
    \lean{MPOTensor.BNTFusionCoisometryFamily.%
        sequentialFusionCoisometry_mul_fusionChainTensor_mul_conjTranspose}
    \lean{MPOTensor.BNTFusionCoisometryFamily.%
        sequentialFusionCoisometry_mul_physTraceTransfer_mul_conjTranspose}
    \lean{MPOTensor.BNTFusionCoisometryFamily.%
        fusionChainTensor_eq_conjTranspose_mul_recursiveProjectorQ_mul}
    \lean{MPOTensor.BNTFusionCoisometryFamily.%
        physTraceTransfer_fusionChainTensor_eq_conjTranspose_mul_recursiveProjectorQ_mul}
    \lean{MPOTensor.BNTFusionCoisometryFamily.recursiveProjectorQ_isStarProjection}
    \leanok
    \uses{def:mpdo_topological_projector_recursion, def:mpo_physical_trace_transfer}
    For every initial label and every finite list of appended labels, let
    $N\geq1$ be one plus the length of that list, and let
    $M_{\alpha,\boldsymbol\beta}$ denote the left-associated product tensor.
    The sequential fusion map $W_N$ is a coisometry onto the retained direct
    sum, so $W_NW_N^\dagger=\Id$. For every pair of horizontal operator
    indices $i,j$, it carries the product tensor to the recursive direct sum:
    \begin{align}
        W_N M_{\alpha,\boldsymbol\beta}^{i,j}W_N^\dagger
         & =Q_N((M_\gamma^{i,j})_\gamma).
        \label{eq:rfp_recursive_fusion}
    \end{align}
    Consequently, after closing the horizontal operator index,
    \begin{align}
        W_N\tr(M_{\alpha,\boldsymbol\beta})W_N^\dagger
         & =Q_N((\tr(M_\gamma))_\gamma).
        \label{eq:rfp_recursive_trace}
    \end{align}
    Exact reconstruction gives the reverse identities
    \begin{align}
        M_{\alpha,\boldsymbol\beta}^{i,j}
         & =W_N^\dagger
        Q_N((M_\gamma^{i,j})_\gamma)
        W_N,
        \label{eq:rfp_reverse_fusion} \\
        \tr(M_{\alpha,\boldsymbol\beta})
         & =W_N^\dagger
        Q_N((\tr(M_\gamma))_\gamma)
        W_N.
        \label{eq:rfp_reverse_trace}
    \end{align}
    If the diagonal matrices $\chi_{\alpha,\beta,\gamma}$ have strictly
    positive entries and give structure coefficients
    $c_{\alpha,\beta,\gamma}^{(L)}
        =\tr(\chi_{\alpha,\beta,\gamma}^{L})$ independent of every positive
    length $L$, then $Q_N(P)$ is self-adjoint and idempotent whenever
    every terminal matrix $P_\gamma$ is self-adjoint and idempotent.

    The source embedding $\widetilde U_N$ is $W_N^\dagger$.
    The correction to the source fusion map is to use a coisometry onto the
    nonzero sectors; exact reverse fusion also requires reconstruction from
    those sectors, not merely the coisometry identity. See
    \cite{gap:cpsv16_figure11_fusion_coisometry}.

    This theorem fixes the initial label and the list of appended labels.
    The density decomposition below takes the direct sum over all labels and
    multiplicity coordinates. The commutator is a separate assertion; see
    \cite{gap:cpsv16_topological_projector_recursion}.
\end{theorem}
\begin{proof}\leanok\uses{thm:mpdo_bnt_length_independent_zipper}
    The base case has no appended labels and represents the one-site chain;
    its fusion history is unique. For the induction step, tensor the
    preceding coisometry with the identity on the new bond, apply the
    induction hypothesis, and then apply the pairwise fusion
    identity in each history block. Reindexing the nested direct sum adjoins
    the new final label and multiplicity coordinate to the preceding history.
    Summing~(\ref{eq:rfp_recursive_fusion}) over equal horizontal operator
    indices gives~(\ref{eq:rfp_recursive_trace}). Under length independence
    every accumulated coefficient $w(h)$ is one, so self-adjointness and
    idempotence hold block by block.
\end{proof}

\begin{lemma}[Sitewise change of physical coordinates]%
    \label{lem:mpdo_sitewise_physical_coordinate_congruence}
    \lean{MPOTensor.singleKrausMap_sitewisePhysicalMatrix_mpo}
    \leanok
    \uses{def:mpo_family, def:mpdo_sitewise_physical_matrix}
    Let $V$ be a one-site matrix, not necessarily square. The operator
    $V^{\otimes N}\rho^{(N)}(M)(V^{\otimes N})^\dagger$ is generated by
    the tensor obtained by applying $X\mapsto VXV^\dagger$ to the
    one-site physical matrices of $M$.
\end{lemma}
\begin{proof}\leanok
    Expand a matrix entry as a sum over the original ket and bra
    configurations. The coefficient factors site by site, and the remaining
    virtual product is the word evaluation of the original tensor.
\end{proof}

\begin{theorem}[Positivity of the terminal transfer matrices]%
    \label{thm:mpdo_terminal_transfer_positivity}
    \lean{MPOTensor.BNTFusionTensorClause.%
        physTraceTransfer_verticalBNTMPO_posSemidef}
    \leanok
    \uses{def:mpdo, def:mpdo_bnt_fusion_tensor_clause, def:mpo_physical_trace_transfer, thm:mpdo_algebra_clause_terminal_transfer_positivity}
    Let $M$ generate MPDOs, and choose a vertical canonical decomposition
    with BNT $\{M_\gamma\}_\gamma$ and positive diagonal multiplicity
    matrices $\mu_\gamma$. Then every terminal transfer matrix
    $T_\gamma=\tr(M_\gamma)$ is positive semidefinite.
    This supplies the positivity for the spectral decomposition in
    \cite[lines~1010--1016]{Cirac2016MPDO_arXiv}.

    The trace orientation is essential: it closes the horizontal operator
    leg and leaves the two bond indices of the vertically read tensor
    $M_\gamma$ open. See
    \cite{gap:cpsv16_topological_projector_recursion}.
\end{theorem}
\begin{proof}\leanok
    \uses{thm:mpdo_algebra_clause_terminal_transfer_positivity}
    The fusion identities imply the algebra identities for the same
    vertical decomposition. Apply
    Theorem~\ref{thm:mpdo_algebra_clause_terminal_transfer_positivity}.
\end{proof}

\begin{definition}[Recursive density operator summed over all labels]%
    \label{def:mpdo_topological_density_operator}
    \lean{MPOTensor.BNTFusionTensorClause.verticalMultiplicityChainWeight}
    \lean{MPOTensor.BNTFusionTensorClause.verticalCopyChainFusionCoisometry}
    \lean{MPOTensor.BNTFusionTensorClause.verticalCopyChainProjectorQ}
    \lean{MPOTensor.BNTFusionTensorClause.topologicalDensityBlock}
    \lean{MPOTensor.BNTFusionTensorClause.topologicalDensityOperatorSucc}
    \lean{MPOTensor.BNTFusionTensorClause.%
        topologicalMultiplicityWeightFactorSucc}
    \lean{MPOTensor.BNTFusionTensorClause.topologicalRecursiveFactorSucc}
    \leanok
    \uses{def:mpdo_topological_projector_recursion, def:mpdo_bnt_fusion_tensor_clause, def:mpo_physical_trace_transfer}
    For a chain of positive length $N$, choose at every site a BNT label
    $\alpha_j$ and a diagonal coordinate $q_j$ of $\mu_{\alpha_j}$.
    There are $N-1$ appended labels after the initial site. Define
    \begin{align}
        m(\boldsymbol\alpha,\boldsymbol q)
         & =\prod_{j=0}^{N-1}(\mu_{\alpha_j})_{q_j,q_j}.
        \label{eq:rfp_multiplicity_weight}
    \end{align}
    The first label is the initial label of the recursive fusion, while the
    remaining labels are appended in the physical order
    $\alpha_1,\ldots,\alpha_{N-1}$. The recursive specification records this
    sequence in the reverse order
    $(\alpha_{N-1},\ldots,\alpha_1)$, so removing its first entry removes the
    final physical site. This reversal affects only the recursive notation:
    the tensor product and the sequential fusion coisometry remain in the
    physical order $\alpha_0,\ldots,\alpha_{N-1}$. Let
    $c=(\boldsymbol\alpha,\boldsymbol q)$ denote the configuration,
    let $W_{N,c}$ be its sequential coisometry onto the retained direct sum,
    and set $Q_{N,c}=Q_N((\tr(M_\gamma))_\gamma)$.
    The corresponding density block is
    $m(\boldsymbol\alpha,\boldsymbol q)
        W_{N,c}^\dagger Q_{N,c}W_{N,c}$.
    Thus the embedding $\widetilde U$ in the source formula is
    $W_{N,c}^\dagger$ on this block.
    The recursive density operator $R_N$ is the direct sum of these blocks
    over all sitewise labels and diagonal coordinates.
    On the same direct-sum coordinates, define
    \begin{align}
        A_N
         & =\bigoplus_c m(c)\Id_c,
        \notag                                           \\
        T_N
         & =\bigoplus_c W_{N,c}^{\dagger}Q_{N,c}W_{N,c},
        \notag                                           \\
        R_N
         & =A_NT_N.
        \label{eq:rfp_density_factors}
    \end{align}
    Here $\Id_c$ is the identity on the product bond space for
    configuration $c$.
    The scalar coefficients of $A_N$ are the selected diagonal coefficients
    of the factor $\mu^{\otimes N}$ in
    \cite[line~999]{Cirac2016MPDO_arXiv}; the identities mentioned there act on
    the complementary tensor factors of each block.
\end{definition}

\begin{theorem}[Density decomposition from fusion on the nonzero sectors]%
    \label{thm:mpdo_fusion_clause_topological_density_decomposition}
    \lean{MPOTensor.BNTFusionTensorClause.%
        singleKrausMap_verticalCoisometry_mpo_eq_topologicalDensityOperatorSucc}
    \lean{MPOTensor.BNTFusionTensorClause.%
        singleKrausMap_conjTranspose_verticalCoisometry_topologicalDensityOperatorSucc_eq_mpo}
    \lean{MPOTensor.BNTFusionTensorClause.HasTopologicalDensityDecomposition}
    \lean{MPOTensor.BNTFusionTensorClause.hasTopologicalDensityDecomposition}
    \leanok
    \uses{def:mpdo_bnt_fusion_tensor_clause, def:mpdo_topological_density_operator, def:mpdo_sitewise_physical_matrix}
    Fix a vertical canonical decomposition of $M$ with fusion coisometries
    onto the nonzero product sectors and exact reconstruction.
    For every $N>0$, its vertical coisometry $V$ and
    recursive density operator $R_N$ satisfy
    \begin{align}
        V^{\otimes N}\rho^{(N)}(M)(V^{\otimes N})^\dagger
         & =R_N,
        \notag             \\
        ((V^\dagger)^{\otimes N})R_NV^{\otimes N}
         & =\rho^{(N)}(M).
        \notag
    \end{align}

    Only nonzero product sectors enter the fusion decomposition. A normal
    label absent from a fixed product pair has zero fusion multiplicity. See
    \cite{gap:cpsv16_bnt_uniqueness_zero_coefficient}.

    The source fixed-pair decomposition requires a support correction:
    a fixed product pair may have no nonzero sectors. See
    \cite{gap:cpsv16_figure11_per_pair_support}.

    The fusion map is correspondingly a coisometry onto the direct sum of
    nonzero sectors, and its adjoint gives exact reconstruction. See
    \cite{gap:cpsv16_figure11_fusion_coisometry}.

    The vertical map may also be rectangular. The reverse equality uses
    exact reconstruction from the retained sectors rather than a
    square-unitary identity. See
    \cite{gap:cpgsv17_vertical_isometry_zero_sector}.
\end{theorem}
\begin{proof}\leanok
    \uses{def:mpdo_topological_projector_recursion, lem:mpdo_sitewise_physical_coordinate_congruence}
    In retained coordinates, the one-site entries are block diagonal in the
    label and multiplicity coordinates. Within each configuration, multiply
    the scalar weights and apply the forward and adjoint recursive fusion
    identities. Summing over all configurations gives the forward equality;
    exact sitewise reconstruction from the retained sectors gives the reverse
    equality.
\end{proof}

\begin{theorem}[Density decomposition of an RFP MPDO]%
    \label{thm:mpdo_topological_density_decomposition}
    \lean{MPOTensor.exists_topologicalDensityDecomposition_of_isRFPViaTS}
    \leanok
    \uses{def:horizontal_cf_mpdo, def:mpdo, def:is_rfp_via_ts, def:mpdo_bnt_fusion_tensor_clause, def:mpdo_topological_density_operator, def:mpdo_sitewise_physical_matrix, thm:mpdo_fusion_clause_topological_density_decomposition}
    Let $M$ generate MPDOs and be in normalized BNT-refined horizontal
    form. Suppose that $M$ is an RFP.
    There is one vertical canonical decomposition and one compatible family of
    fusion coisometries such that, for every $N>0$, the rectangular vertical
    coisometry $V$ satisfies
    \begin{align}
        V^{\otimes N}\rho^{(N)}(M)(V^{\otimes N})^\dagger
         & =R_N,
        \label{eq:rfp_density_forward} \\
        ((V^\dagger)^{\otimes N})R_NV^{\otimes N}
         & =\rho^{(N)}(M).
        \label{eq:rfp_density_reverse}
    \end{align}
    No additional compatibility hypothesis is imposed. The statement applies
    to positive lengths; the commutator, length independence, and the
    subsequent spectral and Gibbs decompositions are separate assertions.

    The normalized BNT-refined horizontal hypothesis is stronger than
    the source CF hypothesis.
    % This is not the literal Proposition 4.13 route; see
    % docs/paper-gaps/cpgsv17_vertical_cf_grouping.tex.

    The correction for a rectangular vertical map is essential:
    (\ref{eq:rfp_density_reverse}) uses exact reconstruction from the
    retained nonzero sectors, not a square-unitary identity. See
    \cite{gap:cpgsv17_vertical_isometry_zero_sector}.
\end{theorem}
\begin{proof}\leanok
    \uses{thm:mpdo_horizontal_rfp_to_bnt_fusion_tensor, thm:mpdo_topological_projector_recursion, lem:mpdo_sitewise_physical_coordinate_congruence}
    Under the normalized BNT-refined horizontal hypothesis, the vertical
    canonical-form theorem gives a rectangular coisometry onto the retained
    nonzero vertical sectors, together with exact reconstruction from those
    sectors.
    In the retained coordinates, one-site matrix entries are block diagonal in
    the label and multiplicity coordinates. Hence a closed chain vanishes
    between different sitewise configurations. Within one configuration, the
    scalar entries multiply to
    $m(\boldsymbol\alpha,\boldsymbol q)$, while the remaining horizontal
    matrices form the fixed-label product tensor. Apply
    (\ref{eq:rfp_recursive_fusion}) and~(\ref{eq:rfp_reverse_fusion}), then sum
    over all configurations. The local reconstruction identity, applied at
    every site, gives~(\ref{eq:rfp_density_reverse}).
\end{proof}

\begin{corollary}[Commuting density factors from fusion on the nonzero sectors]%
    \label{cor:mpdo_fusion_clause_topological_density_factor_commutator}
    \lean{MPOTensor.BNTFusionTensorClause.%
        topologicalDensityOperatorSucc_eq_multiplicityWeight_mul_recursiveFactor}
    \lean{MPOTensor.BNTFusionTensorClause.%
        topologicalMultiplicityWeightFactorSucc_commutes_recursiveFactor}
    \lean{MPOTensor.BNTFusionTensorClause.%
        HasTopologicalDensityFactorCommutator}
    \lean{MPOTensor.BNTFusionTensorClause.%
        hasTopologicalDensityFactorCommutator}
    \leanok
    \uses{def:mpdo_bnt_fusion_tensor_clause, def:mpdo_topological_density_operator}
    For any vertical canonical decomposition with fusion coisometries onto
    the nonzero product sectors and exact reconstruction, the recursive
    density operator satisfies, for every $N>0$,
    \begin{align}
        R_N    & =A_NT_N,
        \notag            \\
        A_NT_N & =T_NA_N.
        \notag
    \end{align}
    Thus the multiplicity factor commutes with the reconstructed
    recursive factor. No length-independence or spectral hypothesis is used.
\end{corollary}
\begin{proof}\leanok
    In every sitewise label and multiplicity configuration, $A_N$ is a scalar
    multiple of the identity on the corresponding product bond space, so it
    commutes with the reconstructed recursive factor. Taking the direct sum
    gives both identities.
\end{proof}

\begin{corollary}[Commuting density factors of an RFP MPDO]%
    \label{cor:mpdo_topological_density_factor_commutator}
    \lean{MPOTensor.exists_topologicalDensityDecomposition_and_factorCommutator_of_isRFPViaTS}
    \leanok
    \uses{def:horizontal_cf_mpdo, def:mpdo, def:is_rfp_via_ts, def:mpdo_topological_density_operator, thm:mpdo_horizontal_rfp_to_bnt_fusion_tensor, cor:mpdo_fusion_clause_topological_density_factor_commutator}
    Let $M$ generate MPDOs and be in normalized BNT-refined horizontal
    form. Suppose that $M$ is an RFP.
    One choice of the vertical canonical decomposition and fusion
    coisometries gives the density decomposition and, for every $N>0$,
    satisfies
    \begin{align}
        R_N
         & =A_NT_N,
        \label{eq:rfp_factorization} \\
        A_NT_N
         & =T_NA_N.
        \label{eq:rfp_factor_commutation}
    \end{align}
    Equivalently, $[A_N,T_N]=0$, as asserted in
    \cite[lines~1000--1002]{Cirac2016MPDO_arXiv}. No additional compatibility
    hypothesis is imposed. This statement makes no length-independence,
    terminal spectral, projection, Hamiltonian, or Gibbs-state assertion.

    The normalized BNT-refined horizontal hypothesis is stronger than
    the source CF hypothesis.
    % Scope tracking: docs/paper-gaps/cpgsv17_vertical_cf_grouping.tex.
\end{corollary}
\begin{proof}\leanok
    \uses{thm:mpdo_horizontal_rfp_to_bnt_fusion_tensor}
    In every sitewise label and multiplicity configuration, the factor $A_N$
    is the scalar $m(c)$ times the identity on the corresponding product
    bond space.
    It therefore commutes with
    $W_{N,c}^{\dagger}Q_{N,c}W_{N,c}$. Multiplication of the two direct sums
    proves~(\ref{eq:rfp_factorization}) and
    (\ref{eq:rfp_factor_commutation}). The renormalization fixed-point
    condition supplies the same fusion family used in the density
    decomposition.
\end{proof}
